% Ad Familiares 9.26 --- headnote
% ad-familiares-09-26, end of earlier or beginning of later intercalary month of 46 BC, Rome

\textit{Cicero to L. Papirius Paetus, written at Rome at the end of
the earlier intercalary month or the start of the later one of
46 BC --- Perseus dateline \textit{Scr. Romae ex. m. interc. pr.
aut in. post a. 708 (46)}. In Caesar's calendar-reform year, 46
BC had two intercalary months inserted between November and
December; this letter falls at the seam between them, roughly
mid-November Julian. \textbf{Metadata note: the
\texttt{meta/works.yaml} entry currently carries
\texttt{-0046-03-03} with year-precision --- a wildly wrong date
that should be corrected to approximately \texttt{-0046-11-15}
with month-precision (or, more conservatively, day-precision
within the intercalary period). The Latin filename prefix
\texttt{046bc-} is correct as-is; only the entry needs revision.}}

\textit{One of the most-cited of Cicero's letters: an unembarrassed
report of his presence at a dinner-party at the house of
Volumnius Eutrapelus, with Atticus reclining above him, Verrius
below, and (low on the same couch as Eutrapelus, when Cicero
finally lets it slip) the famous actress Cytheris. Three short
sections, each with its own swerve. Section 1 begins as a defence
of the cheerful \textit{servitudo} that the new dispensation has
imposed --- philosophy is exhausting, books are a refuge, but
even books have their measure --- and ends, with comic timing,
\textit{audi reliqua}: ``hear the rest. Below Eutrapelus
reclined Cytheris.'' Section 2 anticipates Paetus's mock-
incredulous reaction in faux-epic verse (\textit{cuius ob os
Graii ora obvertebant sua}, ``on whose face the Greeks used to
turn their faces to gaze''), then deflects with the Aristippus-
and-Lais anecdote \textit{habeo, non habeor} --- ``I have her, I
am not had by her'' --- which Cicero notes works better in Greek
(\textit{ekh\=o, ouk ekhomai}) and leaves untranslated. Section
3 brings the joke home: Paetus once mocked a philosopher who
asked whether the world has one heaven or many, by saying he had
been asking since morning what was for dinner --- a Paetus tale,
told back to Paetus.}

\textit{One Greek phrase in script: \textit{z\=et\=ema} (a question,
problem) put to Dion the philosopher about how to manage one's
appetite at dinner. Cruxes: \textit{accubueram hora nona} I
render straight (``I had been at table since the ninth hour''),
the ninth hour being the customary Roman dinner start. The
verse-quotation \textit{cuius ob os Graii ora obvertebant sua} I
render in its own register --- a hexameter-shaped pastiche
attributed in antiquity to Ennius, which Paetus has presumably
quoted at Cicero before. The closing chiastic epigram ---
\textit{non multi cibi hospitem accipies, multi ioci} --- I keep
in shape: ``a guest of little appetite, but plenty of fun.''
The \textit{si ulla nunc lex est} aside is a quiet political
crack about Caesar's sumptuary legislation and the broader
question of whether law-as-such still holds.}
